\section{Cluster structures on braid varieties}\label{sec:braid_variety_cluster}

The object of this section is to explain the construction of a cluster algebra structure on the commutative algebra $\Z[X(\beta)]$ given by the coordinate ring of regular functions on $X(\beta)$. In order to do that, we first review framed flags in \Cref{ssec:framed_rep} and discuss certain quantitative measurements of transversality between them in \Cref{ssec:rational_functions}, which help towards better understanding the tori underlying the cluster seeds for $\Z[X(\beta)]$. The cluster $\SA$-variables are constructed in \Cref{ssec:cluster_Avariables_weave} and \Cref{ssec:cluster_A_variables_propagation_vertices} explains how to obtain them algorithmically by scanning the weave bottom-to-top. Similarly, \Cref{ssec:cluster_Xvariables_weave} constructs the dual cluster $\SX$-variables and presents their key properties.\\

In \Cref{ssec:A_and_X_schemes_weave} we explain the relation between the braid variety $X(\beta)$ and the cluster $\SA$ and $\SX$-schemes, and \Cref{ssec:brick_compactifications} concludes by discussing certain projective compactifications of braid varieties that can be used to better understand them, cf.~\Cref{sssec:applications_compactifications}.\\

\begin{comment}
Theorem \ref{thm:main} proves that the rings of regular functions of each such types of varieties are cluster algebras. Even in cases where Theorem \ref{thm:main} was known, e.g.~double Bruhat cells and positroids, the proof we give -- following \cite{CGGLSS25} -- is different from the original sources. In these cases, some comparisons with the previously known cluster algebras are discussed in \cite[Section 10]{CGGLSS25}, \cite[]{CLSW23}, \cite[Section 10.5]{GLSBS} and \cite{clustercomparison25}.
\end{comment}

%%%%%%%%%%%%%%%%%%%%%%%%%%%%%%%%%%%%%%%%%%%%%
%%%%%%%%%%%%%%%%%%%%%%%%%%%%%%%%%%%%%%%%%%%%%
\subsection{Framed representatives for points in $X(\beta)$}\label{ssec:framed_rep} Let $\beta\in\brn$ be a positive braid word. The braid variety $X(\beta)$ and the weave tori $T_\w\sse X(\beta)$ are both defined by using relative positions between flags in $\G/\borel$. For a pair $\borel_1,\borel_2\in\G/\borel$ we can decide their relative position, given by a permutation in $W(\SL_n)$, but there is not enough information in $\G/\borel$ to measure this relative position in a quantifiable manner. For instance, if $\borel_1\stackrel{s_i}{\to}\borel_2$, we would want to be able to measure how (or how much) $\borel^i_1$ and $\borel^i_2$ differ. This leads to the notion of framed flags. To define them, let $\uni\sse \borel$ be a fixed unipotent radical, e.g.~the subgroup of upper unitriangular matrices for $\G=\SL_n$.

\begin{definition}[Framed flags]
$\G/\uni$ is said to be the principal flag variety, and a framed flag is a point in $\G/\uni$. Given $w\in W(G)$, two framed flags $x\uni,y\uni\in\G/\uni$ are said to be in relative $w$-position, denoted $x\uni\stackrel{w}{\to}y\uni$, if $x^{-1}y\in \uni w\uni$.\hfill$\Box$
\end{definition}
Since $\T\cong\borel/\uni$, the projection map $\pi:\G/\uni\lr\G/\borel$ induced by the inclusion $\uni\sse \borel$ is a principal $\T$-bundle; in particular its fibers are isomorphic to $\T$. In the correspondence between Borel subgroups and flags for $\G=\SL_n$, as in \cref{rmk:borels_flags}, framed flags in $\G/\uni$ can be interpreted as pairs $(\FF,\{f_1,\ldots,f_n\})$ consisting of a full flag
$$\FF=(\FF^0\subset \FF^1\subset \FF^2\subset\ldots \subset \FF^n)\in\G/\borel,\quad \FF^i\cong k^i,$$
equipped with a set of generators $\{f_1,\ldots,f_n\}$ of the 1-dimensional vector spaces $\FF^i/\FF^{i-1}=\langle f_i\rangle$ given by the consecutive quotients of the full flag $\FF$ such that these generators lift to vectors $\tilde{f}_i\in k^n$ with volume one, i.e.~$\tilde{f}_1\wedge\ldots\wedge \tilde{f}_n=\mbox{vol}_n$, where $\mbox{vol}_n$ is the volume element intrinsically defining $\SL_n\sse\on{GL}_n$. Note that the projection $\pi(\FF,\{f_1,\ldots,f_n\})=\FF$ above is forgetting the data of the choice of generators and just remembering the underlying flag $\FF$. Each flag in $\G/\borel$ can be represented by many framed flags, and such framed flags are a $\T$-torsor: this simply corresponds to the choices of generators $\{f_1,\ldots,f_n\}$.

\begin{example}[Framed flags in $\SL_2$]\label{ex:framed_flags_SL2}
Consider $\G=\SL_2$, its framed flag variety is $\SL_2/\uni\cong\kk^2 \setminus \{(0,0)\}$. For instance, if we represent $g\in\SL_2$ via $g=(v_1\quad v_2)$ where $v_1,v_2\in\kk^2$ denote the columns of $g$, then an isomorphism is given by $g=(v_1\quad v_2)\uni\mapsto v_1$. The forgetful projection $\pi:\SL_2/\uni\lr\SL_2/\borel$ is the quotient of $\kk^2 \setminus \{(0,0)\}$ by the $\bG_m$-action diagonally scaling the entries of $v_1\in \kk^2 \setminus \{(0,0)\}$. In particular, the fibers of $\pi$, which parametrize the possible framings of a given flag $g\borel\in\SL_2/\borel$, are the $\bG_m$-orbits and thus isomorphic to $\T=\bG_m$, since the action is free. In the coordinates above, the fiber over $g\borel$ for $g=(v_1\quad v_2)$ consists of those framed flags $\tilde{g}\uni$ such that
\begin{equation}\label{eq:framed_flags_sl2}
\tilde{g}=(av_1\quad a^{-1}v_2+bv_1),\quad \mbox{for some }(a,b)\in\kk^\times\kk.
\end{equation}
Varying $a\in\kk^\times$ provides the $\G_m$-fiber, while varying $b$ does not change the unipotent coset. As in \cref{rmk:borels_flags},  the flag $g\borel$ is represented by the sequence of nested subspaces $\{0\}\subset \langle v_1\rangle \subset \kk^2\cong\langle v_1,v_2\rangle$. Since all framed flags $g'\uni$ with $\pi(g'\uni)=g\borel$ are parametrized by \eqref{eq:framed_flags_sl2}, the framing $\{f_1,f_2\}$ of such a framed flag $g'\uni$ is given by the image $f_1=av_1$ of $v_1$ in $\langle v_1\rangle/\{0\}$ and $f_2=a^{-1}v_2$ in $\kk^2/\langle v_1\rangle$. Note that $av_1\wedge(a^{-1}v_2+\lambda v_1)=v_1\wedge v_2=\on{vol}_{\kk^2}$ since $v_1,v_2$ are the columns of a matrix in $\SL_2$. In general, given a framed flag $g'\uni\in\SL_n/\uni$, the framing $f_i$ at the $i$th quotient of the flag $\pi(g'\uni)$ is given by the image of the $i$th column of $g'$.\sbs
\end{example}

For us, it is useful that {\it every} point $(\borel_1,\borel_2,\ldots,\borel_{\ell+1})\in X(\beta)$ in the braid variety can be represented by a tuple of framed flags and, moreover, uniquely so if we require the framed lift of $\borel_1$ to be the specific lift $\uni$ of $\borel_1=\borel$. To show that, it suffices to note that given any framed flag $x\uni\in\G/\uni$ and (unframed) flag $\FF\in\G/\borel$ with $x\borel\stackrel{s_i}{\to}\FF$, there exists a {\it unique} framed flag $y\uni\in\G/\uni$ such that $x\uni\stackrel{s_i}{\to}y\uni$. Indeed, at an $s_i$-crossing with $x\borel$ on its left and $\FF$ on its right, there are canonical isomorphisms
$$\frac{(x\borel)^k}{(x\borel)^{k-1}}\cong \frac{\FF^k}{\FF^{k-1}}\mbox{ if }k\neq i,i+1,\quad \frac{(x\borel)^i}{(x\borel)^{i-1}}\cong \frac{\FF^{i+1}}{\FF^{i}}\mbox{ and }\quad \frac{(x\borel)^{i+1}}{(x\borel)^{i}}\cong \frac{\FF^{i}}{\FF^{i-1}}$$
and thus a framing of $x\borel$, lifting it to a framed flag, is uniquely transferred to a framing for $\FF$. Therefore, given $(\borel_1,\borel_2,\ldots,\borel_{\ell+1})\in X(\beta)$ we can consider the {\it unique} tuple
$$(\uni_1,\uni_2,\ldots,\uni_{\ell+1}),\quad \uni_i\in\G/\uni,\quad \mbox{s.t. }\uni_j\stackrel{s_{i_j}}{\uni_{j+1}}\mbox{ and }\uni_1=\uni.$$
\noindent By construction, $(\pi(\uni_1),\pi(\uni_2),\ldots,\pi(\uni_{\ell+1}))\in X(\beta)$.

\begin{remark}\label{rmk:unipotent_representative_lastflag}
This lifting procedure frames each flag in $(\borel_1,\borel_2,\ldots,\borel_{\ell+1})\in X(\beta)$ by first framing $\borel$ by $\uni$ and propagating the framing to the right. For instance if $\delta(\beta)=w_0$, this process frames the last flag $\borel_{\ell+1}=w_0\borel$, but the framing is {\it not} necessarily $w_0\uni$. Rather, the last framing is given by $tw\delta(\beta)\uni$ for some element $t\in\T$, but not necessarily the identity. This explains geometrically the map $\mu:X(\beta)\lr\T$ in \Cref{rmk:torus_quasimoment_map}.\sbs
\end{remark}

%%%%%%%%%%%%%%%%%%%%%%%%%%%%%%%%%%%%%%%%%%%%%
%%%%%%%%%%%%%%%%%%%%%%%%%%%%%%%%%%%%%%%%%%%%%
\subsection{Rational functions on $X(\beta)$ via weave edges}\label{ssec:rational_functions} Let $\beta\in\brn$ be a positive braid, $X(\beta)$ its braid variety, $\w:\beta\lr\delta(\beta)$ a Demazure weave and $T_\w\sse X(\beta)$ its weave torus. The Demazure weave $\w$ not only provides $T_\w$, it also provides a set of coordinates on $T_\w$. In order to explicitly describe them, we fix a pinning $(x_{i}, y_{i})_{i \in \dynkin}$ of the group $\G$, i.e.~ for every simple root indexed by $i \in \dynkin$ in its Dynkin diagram we choose two isomorphisms $x_{i}: \kk \to \uni_{i}^{+}$ and $y_{i}: \kk \to \uni_{i}^{-}$, where $\uni_{i}^{+}$ and $\uni_{i}^{-}$ are the corresponding root subgroups of $i \in \dynkin$, such that the assignment
\[
\left(\begin{matrix} 1 & z \\ 0 & 1 \end{matrix}\right) \mapsto x_{i}(z), \quad \left(\begin{matrix} 1 & 0 \\ z & 1 \end{matrix}\right) \mapsto y_{i}(z),\quad \left(\begin{matrix} t & 0 \\ 0 & t^{-1} \end{matrix}\right) \mapsto \chi_{i}(t)
\]
gives an algebraic group morphism $\varphi_{i}: \SL_{2}(k) \to \G$, with $\chi_{i}: \bG_m \to T$ the simple coroot corresponding to $i \in \dynkin$.\footnote{Recall that a simple algebraic group $\G$ admits a pinning and any two pinnings are conjugate.} Fixing a pinning is useful in our context because it allows us to explicitly measure flag transversality and framing differences in a quantifiable manner. Specifically, we have the following result:

\begin{lemma}\label{lem:framing_variables}
Let $g\uni,g'\uni\in\G/\uni$ be a pair of framed flags whose underlying flags are in $s_i$-relative position. If a representative $g$ of $g\uni$ is chosen, then there exist unique values $(\textcolor{purple}{\tilde{z}},\textcolor{blue}{u})\in\kk\times\kk^\times$ such that
$$g'\uni=gx_{i}(\textcolor{purple}{\tilde{z}})s_{i}\chi_i(\textcolor{blue}{u})\uni$$
where $s_{i} := \varphi_{i}\left(\begin{matrix}0 & -1 \\ 1 & 0 \end{matrix} \right)\in \G$, so $s_{i} \in N_{\G}(T)$ lifts the simple reflection corresponding to $i \in \dynkin$.\qed
\end{lemma}

\noindent The importance of \Cref{lem:framing_variables} is that it provides us with explicit values $\textcolor{purple}{\tilde{z}},\textcolor{blue}{u}$ that measure the transversality of two framed flags. That is, given a representative for a framed flag $g\uni$, the variable $\textcolor{purple}{\tilde{z}}$ parametrizes the choices of (unframed) flags which are $s_i$-transverse to the underlying flag $g\borel$, and the variable $\textcolor{blue}{u}$ parametrizes the choices of framing, cf.~\Cref{ex:framed_flags_SL2}.\\

Let us fix a representation of the points of $X(\beta)$ as tuples of framed flags, as in \Cref{ssec:framed_rep}, by choosing the constant functions $u_e=1$ at the top edges of the weave. Now we can use \Cref{lem:framing_variables} to construct more interesting rational functions, as follows. For any point in a weave edge $e\in\w$ labeled by $s_i$, we can consider a horizontal slice $h$ through that point. Given a point $(\borel_1,\ldots,\borel_{\ell_h+1})\in X(\beta_h)$ in the corresponding flag variety, the two flags $\borel_j,\borel_{j+1}$ to the left and right of the weave edge $e$ must be in $s_{i_j}$-relative position, where $\s_{i_j}$ is the $j$th crossing of $\beta_h$ reading left-to-right. Since points $X(\beta)$ are represented as tuples of framed flags, we can use the raking procedure in \cite[Section 5]{CGGLSS25} and obtain a morphism $X(\beta_h)\lr X(\beta)$, which is oftentimes a subvariety, and in this case $(\borel_1,\ldots,\borel_{\ell_h+1})\in X(\beta_h)$ is already lifted to a tuple of framed flags $(\uni_1,\ldots,\uni_{\ell_h+1})\in X(\beta_h)$. By writing $\uni_j=g\uni$ and $\uni_{j+1}=g'\uni$, and choosing a representative $g\in\G$, \Cref{lem:framing_variables} implies that 
there exist unique values $(\textcolor{purple}{\tilde{z_e}},\textcolor{blue}{u_e})\in\kk\times\kk^\times$ such that
$$g'\uni=gx_{i_j}(\textcolor{purple}{\tilde{z_e}})s_{i_j}\chi_{i_j}(\textcolor{blue}{u_e})\uni$$
The value of $\textcolor{blue}{u_e}$, as $(\uni_1,\ldots,\uni_{\ell_h})$ ranges over the allowed tuples of flags in $T_\w$, defines a rational function on the open torus $T_\w$. Thus, by implicitly including the raking corrections at the virtual vertices so as to always compare to $X(\beta)$, as in \cite[Section 5]{CGGLSS25}, we can write:

\begin{definition}[Rational function of a weave edge]\label{def:rational_function_weave_edge}
Let $\beta\in\brn$ be a positive braid word, $\w:\beta\lr\delta(\beta)$ a Demazure weave and $e\in E(\w)$ an edge labeled by $s_i$. By definition, and in the notation above, the rational function $\textcolor{blue}{u_e}:X(\beta)\lr\kk$ associated to $e$ is the unique rational function on $X(\beta)$ such that
\begin{equation}
g'\uni=gx_{i_j}(\tilde{z_e})s_{i_j}\chi_{i_j}(\textcolor{blue}{u_e})\uni
\end{equation}
for all pairs of flags $g\uni,g'\uni$ which appear left and right of $e$ in a configuration of flags that belongs to $T_\w$, where a representative $g$ of $g\uni$ is fixed, for some varying value of $\tilde{z_e}$.\sbs
\end{definition}

\noindent The rational functions in \Cref{def:rational_function_weave_edge} are typically not regular functions on $X(\beta)$, i.e.~they do not belong to $\kk[X(\beta)]$. That said, they already provide a coordinate system in $T_\w$:

\begin{lemma}[Coordinates in $T_\w$ from weave vertices]\label{lem:rational_functions_from_weave_vertices}
Let $\beta\in\brn$ be a positive braid and $\w:\beta\lr\delta(\beta)$ a Demazure weave with $d$ trivalent vertices. Consider the set of edges $e_1,\ldots,e_d$ with $e_i$ the south edge of the $i$th trivalent vertex of $\w$ and write $u_i:=u_{e_i}$. Then the set of rational functions $\{u_1,\ldots,u_d\}\in\kk(X(\beta))$ satisfies
$$T_\w\cong \Spec\kk[u_1^{\pm1},\ldots,u_d^{\pm1}].$$
That is, $u_1,\ldots,u_d$ is a coordinate system on the weave torus $T_\w$.
\end{lemma}

\begin{remark}
In \Cref{def:rational_function_weave_edge} and \Cref{lem:rational_functions_from_weave_vertices} there is a subtlety with regards to $u_e$ being well-defined on an entire edge $e$, rather than depending on the chosen point in $e$. Technically, $u_e$ is only well-defined on the edges of the refined framed weave. This is related to the discussion on raking from \cite[Section 5]{CGGLSS25}, as is the choice of $g$ representing $g\uni$, see specifically \cite[Definition 5.8.(5)]{CGGLSS25}. As said above, we treat this implicitly in these notes and refer to \cite{CGGLSS25} for details.\sbs
\end{remark}

%%%%%%%%%%%%%%%%%%%%%%%%%%%%%%%%%%%%%%%%%%%%%
%%%%%%%%%%%%%%%%%%%%%%%%%%%%%%%%%%%%%%%%%%%%%
\subsection{The cluster $A$-variables of a Demazure weave}\label{ssec:cluster_Avariables_weave} The object of this subsection is to define regular functions $A^\w_v\in\kk[X(\beta)]$ associated to the trivalent vertices of any Demazure weave $\w:\beta\lr\delta(\beta)$ via the dual Lusztig cycles introduced in \Cref{def:dual_Lusztig_cycles}. For that, let $\sigma$ be a trivalent or hexavalent vertex of $\w$ and denote by $e_{c}(\sigma)$, the weave edge incident to $\sigma$ in the cardinal directions $c\in\{\mathsf{nw},\mathsf{ne},\mathsf{s}\}$ of northwest, northeast and south respectively. Similarly denote by $u_{c}(\sigma)$ the rational function associated to $e_{c}(\sigma)$ via \Cref{def:rational_function_weave_edge}.

\begin{fdef}[Cluster $A$-variables of $\w$]\label{def:cluster_Avariables_weave}
Let $\beta\in\brn$ be a positive braid word, $\w:\beta\to\delta(\beta)$ a Demazure weave, $v\in\w$ a trivalent vertex and $\rho_v$ its dual Lusztig cycle. By definition, the rational function $A^\w_v:X(\beta)\lr\PP_k^1$ associated to $v$ is the product
\begin{equation}\label{eq:cluster_A_variables}
\boxed{
A^\w_v:=\prod_{\sigma\in \w_3\cup \w_6}L_\sigma(\rho_v).
}
\end{equation}
where we have denoted by $\w_3$ and $\w_6$ the sets of 3-valent and 6-valent vertices of $\w$ and each factor $L_\sigma(\rho_v)$ is given by
$$L_\sigma(\rho_v):=
\begin{cases}
\displaystyle
\left(\frac{u_{\rm s}}{u_{\rm nw}}\right)^{\rho_v(e_{\rm ne})-\rho_v(e_{\rm nw})},
&\mbox{ if }\rho_v(e_{\rm ne})\ge \rho_v(e_{\rm nw}),\\[1.2em]
\displaystyle
\left(\frac{u_{\rm ne}}{u_{\rm s}}\right)^{\rho_v(e_{\rm nw})-\rho_v(e_{\rm ne})},
&\mbox{ if }\rho_v(e_{\rm nw})\ge \rho_v(e_{\rm ne}).
\end{cases}$$
where we abbreviated $e_{\rm c}:=e_{\rm c}(\sigma)$ and $u_{c}:=u_c(\sigma)$ for $c\in\{\mathsf{nw},\mathsf{ne},\mathsf{s}\}$.\sbs
\end{fdef}
\vspace{0.2cm}
\Cref{eq:cluster_A_variables} makes it apparent that $A^\w_v$ is constructed by using solely the information of the dual Lusztig cycles $\rho_v$ and the rational functions $u_e$ associated to the weave edges $e\in\w$. If $\w$ is implicitly understood by context or fixed, we often abbreviate $A_v:=A^\w_v$. Given a Demazure weave $\w$, it can be useful to notice that it is possible to write each such $u_e$ in terms of the collection $\{A_v\}$ and the Lusztig cycles $\g_v$ of $\w$. Indeed, it follows from \cite[Theorem 5.12]{CGGLSS25} that
\begin{equation}\label{eq:u_via_cluster_A_variables}
\boxed{
u_e=\prod_{v\in\w_3}A_v^{\g_v(e)}.
}
\end{equation}

\begin{remark}\label{rmk:Avars_different_defs} The presentation of the cluster $A$-variables $\{A_v\}$ in \Cref{def:cluster_Avariables_weave} differs from the original presentation $A_v$ in \cite[Section 5]{CGGLSS25}. They can be shown to be equivalent.\sbs
\end{remark}


\Cref{eq:cluster_A_variables} does not make it apparent that $A_v$ is actually a regular function on $X(\beta)$. It is also important to understand their behavior under weave equivalence and weave mutation. The following facts are proven in \cite[Section 5.2 \& 5.3]{CGGLSS25}:

\begin{fthm}\label{thm:cluster_variables_weave}
Let $\beta\in\brn$ be a positive braid word, $X(\beta)$ its braid variety and $\w:\beta\lr\delta(\beta)$ a Demazure weave with $d$ trivalent vertices $v_1,\ldots,v_d$. Then the following hold:

\begin{enumerate}
    \item For every trivalent vertex $v\in\w$, the rational function $A_v\in\kk(X(\beta))$ is a regular function. That is, $A_v\in\kk[X(\beta)]$, so $A_v:X(\beta)\lr\mathbb{A}^1_\kk$ has no poles.\\

    \item For every trivalent vertex $v\in\w$, the rational function $A_v\in\kk[X(\beta)]$ is non-vanishing if and only if at least one of the weights of the Lusztig cycle $\g_v$ at the south boundary of $\w$ does not vanish.\\

    \item The restrictions of $A_{v_1},\ldots,A_{v_d}$ to the weave torus $T_\w\sse X(\beta)$ are a coordinate system, i.e.~
    $$T_\w\cong \Spec\kk[A_{v_1}^{\pm1},\ldots,A_{v_d}^{\pm1}].$$
    Thus $T_\w\sse X(\beta)$ is cut-out as an open subvariety as the non-vanishing of all the regular functions $A_{v_i}$.\\

    \item If $\w_1\simeq\w_2$ is a weave equivalence, the two sets of regular functions $\{A^{\w_1}_{v_1},\ldots,A^{\w_1}_{v_d}\}$ and $\{A^{\w_2}_{v_1},\ldots,A^{\w_2}_{v_d}\}$ coincide.\\

    \item If $\w_1\stackrel{\mu}{\sim}\w_2$ is a weave mutation at the Lusztig cycle $\gamma_{v_k}$ of $\w_1$, then $A^{\w_1}_{v_i}=A^{\w_2}_{v_i}$ if $i\neq k$ and
    \begin{equation}\label{eq:cluster_mutation_weave_variables}
        \displaystyle A^{\w_1}_{v_k}A^{\w_2}_{v_k}=\prod_{\langle \g_i,\g_k\rangle>0} (A^{\w_1}_{v_i})^{\langle \g_i,\g_k\rangle}+\prod_{\langle \g_i,\g_k\rangle<0} (A^{\w_1}_{v_i})^{-\langle \g_i,\g_k\rangle}.
    \end{equation}
\end{enumerate}
\end{fthm}
\vspace{0.2cm}
The identity \eqref{eq:cluster_mutation_weave_variables} in Theorem \ref{thm:cluster_variables_weave}.(5) is known as the cluster exchange relation, cf.~\cite[Definition 3.1.3]{FWZ16}. In particular, if we denote by $\SA(Q)$ the commutative cluster algebra associated to a quiver $Q$, cf.~\cite[Def.~3.1.9]{FWZ16}, we obtain the inclusion
\begin{equation}\label{eq:inclusion_cluster_algebra}
\SA(Q_\w)\sse \kk[X(\beta)]
\end{equation}
where $Q_\w$ is the intersection quiver of Lusztig cycles of any Demazure weave $\w:\beta\to\delta(\beta)$. Specifically, any Demazure weave $\w:\beta\to\delta(\beta)$ provides a cluster seed for $\SA(Q_\w)$ as a subalgebra of $\kk[X(\beta)]$: the cluster seed is given by the torus $T_\w\sse X(\beta)$ endowed with the quiver $Q_\w$ and the set of regular functions $\{A^\w_v\}$, indexed by the trivalent vertices $v\in\w_3$ of $\w$. We will soon see that \eqref{eq:inclusion_cluster_algebra} is an equality.

%%%%%%%%%%%%%%%%%%%%%%%%%%%%%%%%%%%%%%%%%%%%%
%%%%%%%%%%%%%%%%%%%%%%%%%%%%%%%%%%%%%%%%%%%%%
%%%%%%%%%%%%%%%%%%%%%%%%%%%%%%%%%%%%%%%%%%%%%
\subsection{Propagation of weave cluster $A$-variables after adding weave vertex on top}\label{ssec:cluster_A_variables_propagation_vertices} Let $\beta\in\brn$ be a positive braid word of length $\ell=\ell(\beta)$, and $\w:\beta\lr\delta(\beta)$ a Demazure weave. In \Cref{def:cluster_Avariables_weave} we constructed the cluster $A$-variables $A_v^\w\in\kk[X(\beta)]$ in terms of the dual Lusztig cycles $\{\rho_v\}$ of the weave $\w$. Our task now is to analyze how $\{A_v^\w\}$ changes as we modify the braid word $\beta$ according to a braid move and $\s_i^2\lr\s_i$, i.e.~using a weave vertex, cf \Cref{fig:weaves_types_vertices}.\\

For that, let $\w^+$ be a weave obtained from $\w$ by adding exactly one weave vertex $p\in(\w^+\setminus\w)$ to the north boundary of $\w$. Let $h_\pm$ be horizontal slices exactly above and below this new vertex $p$, cf.~\Cref{fig:BraidVariety_DiagramMoves2} and \Cref{fig:BraidVariety_Trivalent2}, and $\beta_\pm$ the corresponding braid words. By construction, $\beta_-=\beta$ and so we denote by $\{A_v\}$ the cluster variables of $\w$, as they are elements of $\kk[X(\beta)]$. The cluster variables of $\w^+$ are denoted by $\{A_v^+\}$ and, if $p$ is a trivalent vertex the additional cluster variable is denoted by $A^+_p$. This notation is consistent with the fact that $\w^+$ coincides with $\w$ below $p$, and thus the trivalent vertices of $\w$ and $\w^+$ below $p$ are identical, and can be indistinguishably referred to by using the letter $v$. We then have a relation between $X(\beta)=X(\beta_-)$ and $X(\beta_+)$ depending on the type of weave vertex $p$:

\begin{enumerate}
\item {\bf Trivalent.} If $p$ is trivalent, then \Cref{lem:braidvarieties_trivalent}, or see \Cref{fig:BraidVariety_Trivalent2} (right), provides an embedding
$$\iota:X(\beta)\times\bG_m\hookrightarrow X(\beta_+).$$
\Cref{prop:Avariables_up_new_trivalent} below describes how to relate $\{A_v\}$ to the restriction of $\{A^+_v\}$ and $A^+_p$ to the image $\mbox{im}(\iota)$.\\

\item {\bf Hexavalent.} If $p$ is hexavalent, then \Cref{lem:braid_invariance}.(1) provides an isomorphism
$$R3:X(\beta)\stackrel{\cong}{\longrightarrow} X(\beta_+).$$
It is a cluster isomorphism but, importantly, it is {\it not} the case that the pull-back of $\{A_v^+\}$ under $R3$ yields the same initial cluster seed from $\w$. Rather, the pull-back of $(T_{\w^+},Q_{\w^+},\{A_v^+\})$ under $R3$ gives a cluster seed in $\kk[X(\beta)]$ obtained by applying certain mutations to the cluster seed $(T_\w,Q_\w,\{A_v\})$ associated to $\w$, cf.~\cite[Remark 2.7]{casals2025comparingclusteralgebrasbraid}. The local model for this sequence of mutations is already present in the theory of wiring diagrams, see e.g.~\cite[Prop.~3.2.3]{FWZ16}.\\

\item {\bf Tetravalent.} If $p$ is tetravalent, then \Cref{lem:braid_invariance}.(2) provides an isomorphism
$$C:X(\beta)\stackrel{\cong}{\longrightarrow} X(\beta_+).$$
It is a cluster isomorphism that satisfies $C^*A_v^+=A_v$.\\
\end{enumerate}

\noindent In conclusion, the genuinely new phenomenon is at a trivalent vertex, so let us now assume $p\in\w^+$ is a trivalent vertex. By \Cref{lem:braidvarieties_trivalent} and \Cref{def:cluster_Avariables_weave}, the image of $\iota:X(\beta)\times\bG_m\hookrightarrow X(\beta_+)$ coincides with
$$\mbox{im}(\iota)=\{A^+_p\neq0\}$$
as subvarieties of $X(\beta_+)$, with the second factor identified as $\bG_m=\Spec[(A_p^+)^{\pm1}]$. Therefore, we can consider the projection $\pi:X(\beta)\times\bG_m\lr X(\beta)$ onto the first factor and ask how $\iota^*A^+_v$ relate to $\pi^*A_v$, as a comparison of regular functions on $X(\beta)\times\bG_m$.\\

\noindent For specificity, let us write $\beta=\beta_1\s_{i_k}\beta_2$ and $\beta_+=\beta_1\s_{i_k}\s_{i_k}\beta_2$, so that the south edge of the trivalent vertex $p$ corresponds to the $k$th crossing $\s_{i_k}$ of $\beta$. In line with above, we denote the trivalent vertices of $\w$ with the letter $v$, and use $p$ for the top trivalent vertex of $\w^+$, which is the only new trivalent of $\w^+$ (not in $\w$). Denote the north boundary edges of $\w$, corresponding to crossings of $\beta_1\s_{i_k}\beta_2$, by $e_1^{top},\ldots,e_\ell^{top}$. By definition, we let $\gamma_p|_\w$ be the restriction of the Lusztig cycle of $p$ in $\w^+$ to the subweave $\w$, i.e.~ $\gamma_p|_\w$ is the unique weave cycle on $\w$ which has weights $\gamma_p|_\w(e_i^{top})=\delta_{ik}$ and otherwise satisfies the tropical Lusztig rules at all the vertices of $\w$. Let us refer to the cluster variables of $\w^+$ associated to trivalent vertices $v$ also in $\w$ as $\{A_v^+\}$.\\

Finally, let $\{\alpha_i\}$ and $\{\alpha_i^\vee\}$ be the set of simple roots and simple coroots of $\G$, indexed by the vertices in $\dynkin$. Given the braid word $\beta=\sigma_{i_1}\cdots\sigma_{i_\ell}$, we need the sequences of root and coroots from \cite[Equation (16)]{CGGLSS25}, namely:
\begin{equation}\label{eq:roots_coroots}
 \vartheta_j:=s_{i_1}\cdots s_{i_{j-1}}(\alpha_{i_j}),
 \qquad
 \vartheta_j^\vee:=s_{i_1}\cdots s_{i_{j-1}}(\alpha_{i_j}^\vee).
\end{equation}


\begin{proposition}\label{prop:Avariables_up_new_trivalent} In the notation above, the cluster $A$-variables of $\w$ and $\w^+$ satisfy
\begin{equation}\label{eq:Avariables_up_new_trivalent}
\iota^*A^+_v=(\iota^*A_p^+)^{c_v}\pi^*A_v,
\end{equation}
\noindent where the exponent in \eqref{eq:Avariables_up_new_trivalent} is
\begin{equation}\label{eq:c_exponent}
 \boxed{
 c_v:=\rho_v(e_k^{top})
 +\langle\rho_v,\gamma_p|_\w\rangle
 +\sum_{k<j}(\vartheta_j,\vartheta_{k}^\vee)\rho_v(e_j^{top})}
\end{equation}
\sbs
\end{proposition}

\noindent The following clarifications on \Cref{prop:Avariables_up_new_trivalent} might be helpful:\\

\begin{enumerate}
    \item The exponent $c_v$ might be negative, zero, or positive. This is not an issue as $A_p^+$, which is typically not a unit in $\kk[X(\beta_+)]$, restricts to a unit in $\mbox{im}(\iota)$ which is where the identity \eqref{eq:Avariables_up_new_trivalent} takes place.\footnote{In particular, it is not a contradiction with the fact that $A_v^+$ is an irreducible function on $X(\beta_+)$.}\\

    \item The expression $\langle\rho_v,\gamma_p|_\w\rangle$ in \eqref{eq:c_exponent} refers to the intersection number as introduced in \Cref{def:intersection_number}, without any boundary intersection terms.\\

    \item The sum
    $$\sum_{k<j}(\vartheta_j,\vartheta_{k}^\vee)\rho_v(e_j^{top})$$
    in \eqref{eq:c_exponent} appears because of the raking procedure in \cite[Section 5.1]{CGGLSS25}, which also explains why only contributions to the right of $p$ are relevant. It can be conceptually understood as a boundary intersection contribution, cf.~\Cref{rmk:boundary_intersections} and \cite[Definition 4.24]{CGGLSS25}.\\
    
    \item The index $k$ in \eqref{eq:c_exponent} is fixed and independent of $v$. It is fixed by the location of the vertex $p$, as we have written $\beta=\beta_1\s_{i_k}\beta_2$ and $\beta_+=\beta_1\s_{i_k}^2\beta_2$ above.
\end{enumerate}

%%%%%%%%%%%%%%%%%%%%%%%%%%%%%%%%%%%%%%%%%%%%%
%%%%%%%%%%%%%%%%%%%%%%%%%%%%%%%%%%%%%%%%%%%%%
\subsection{The cluster $X$-variables of a Demazure weave}\label{ssec:cluster_Xvariables_weave} The object of this subsection is to define rational functions $X^\w_v\in\kk(X(\beta))$ associated to the trivalent vertices of any Demazure weave $\w:\beta\lr\delta(\beta)$ via the Lusztig cycles introduced in \Cref{def:Lusztig_cycles}, as a dual version of \Cref{ssec:cluster_Avariables_weave}. For that, let us denote by $e_1^{b},\ldots,e_r^{b}$ the south edges of $\w:\beta\to\delta(\beta)$, where $r=\ell(\delta(\beta))$ and consider the sequence of roots $\vartheta_j$ in \eqref{eq:roots_coroots} but associated to the permutation $\delta(\beta)$. Given a Lusztig cycle $\g_v$ of $\w$, we consider the following rational function on the $u$-variables:

\begin{equation}\label{eq:boundary_term_Xvariables}
B_\dd(\g_v):=
\displaystyle
\prod_{j=1}^{\ell(\delta(\beta))}u(e_j^{b})^{-\g_v(e_j^{b})-R_v^{(j)}},\qquad\mbox{where }R_v^{(j)}:=\displaystyle\sum_{j<k}\g_v(e_k^{b})(\vartheta_k,\vartheta_j).
\end{equation}
\noindent The term \eqref{eq:boundary_term_Xvariables} is a boundary contribution, in line with \Cref{rmk:boundary_intersections}. It is only possibly non-trivial for frozen Lusztig cycles, that is, $B_\dd(\g_v)=1$ for any mutable Lusztig cycle $\g_v$.\\

\begin{fdef}[Cluster $X$-variables of $\w$]\label{def:cluster_Xvariables_weave}
Let $\beta\in\brn$ be a positive braid word, $\w:\beta\to\delta(\beta)$ a Demazure weave, $v\in\w$ a trivalent vertex and $\g_v$ its Lusztig cycle. By definition, the rational function $X^\w_v:X(\beta)\lr\PP_k^1$ associated to $v$ is the product
\begin{equation}\label{eq:cluster_X_variables}
\boxed{
X^\w_v:=B_\dd(\g_v)\prod_{\sigma\in \w_3\cup \w_6}\mathcal{L}_\sigma(\g_v).
}
\end{equation}
where we have denoted by $\w_3$ and $\w_6$ the sets of 3-valent and 6-valent vertices of $\w$ and each factor $\mathcal{L}_\sigma(\g_v)$ is given by
$$\mathcal{L}_\sigma(\g_v):=
\displaystyle
u_{\rm nw}^{\g_v(e_{\rm ne})-\g_v(e_{\rm s})}u_{\rm s}^{\g_v(e_{\rm nw})-\g_v(e_{\rm ne})}u_{\rm ne}^{\g_v(e_{\rm s})-\g_v(e_{\rm nw})},$$
where we abbreviated $e_{\rm c}:=e_{\rm c}(\sigma)$ and $u_{c}:=u_c(\sigma)$ for $c\in\{\mathsf{nw},\mathsf{ne},\mathsf{s}\}$.\sbs
\end{fdef}
\vspace{0.2cm}

In order to relate the cluster $A$-variables from \Cref{def:cluster_Avariables_weave} to the cluster $X$-variables in \Cref{def:cluster_Xvariables_weave}, the exchange matrix $B=(B_{ij})$ associated to $Q_\w$ can be modified to the $p$-map, as in \cite[Section 8.2]{CGGLSS25}, between the cluster $A$-variety and the cluster $X$-variety. In short, we can construct a matrix $(p_{ij})$ such that:

\begin{enumerate}
\item The entries $B_{ij}=p_{ij}$ coincide if one of the vertices associated to $i$ and $j$ are mutable.
    \item The entries $p_{ij}\in\Z$ are integers.
    \item $(p_{ij})$ is unimodular, i.e.~$\det((p_{ij}))=\pm1$.
\end{enumerate}

\noindent Note nevertheless that $p$ might have non-zero diagonal entries due to the frozen vertices. We refer to \cite[Section 8.2]{CGGLSS25} for the precise expression of the correct $(p_{ij})$ as $B_{ij}$ plus a (half) correction summand. \\

\begin{lemma}\label{lem:X_vars_via_A_vars} Let $\w:\beta\to\delta(\beta)$ be a Demazure weave and $v_i\in\w$ its $i$th trivalent vertex. Then
\begin{equation}
\label{eq:X_vars_via_A_vars}
\boxed{\pi^*X^\w_{v_i}= \prod_{v_j\in\w_3} (A^\w_{v_j}) ^{p_{ij}}}
\end{equation}
where $(p_{ij})$ is a unimodular integral matrix with $p_{ij}=\langle\g_i,\g_j\rangle$ if $\g_i$ or $\g_j$ are mutable Lusztig cycles.\sbs
\end{lemma}

\noindent At this stage, we read \eqref{eq:X_vars_via_A_vars} as an equality of rational functions on $X(\beta)$ but see \Cref{ex:X_in_terms_of_A} below. The analogue of \Cref{ssec:cluster_A_variables_propagation_vertices} and \Cref{prop:Avariables_up_new_trivalent} for cluster $X$-variables is simpler. The behavior under adding a hexavalent and tetravalent vertex on top, as in \Cref{ssec:cluster_A_variables_propagation_vertices}.(2) and (3) is identical to that of the cluster $A$-variables. In contrast, since cluster $X$-variables are constructed from Lusztig cycles $\{\g_v\}$, rather than dual Lusztig cycles $\{\rho_v\}$, the cluster $X$-variables under adding a trivalent vertex on top do not significantly change. Specifically, the analogue of \Cref{eq:Avariables_up_new_trivalent} for the cluster $X$-variables is
\begin{equation}\label{eq:Xvariables_up_new_trivalent}
 \boxed{\iota^*X^+_v=X_v}\quad\mbox{if }v\neq p,
\end{equation}
and one can express the pull-back under $\iota$ of the unique new cluster $X$-variable $X^+_p$ in $\w^+$ as
$$\boxed{\iota^*X^+_p=(\iota^*A_p^+)^{-1}\prod_{v\in\w_3}(\pi^*X_v)^{q_v}}$$
for some exponents $q_v$ that depend on the extended $p$-map $(p_{ij})$. Namely, if we extend the $p$-map $(p_{ij})$ of $\w$ to a $p$-map for $\w^+$ by adding\footnote{Let us use the convention here that the new vertex is added with the last index.} a row to the bottom of $(p_{ij})$ and a column to its right, then $q=(q_v)$ is the vector obtained by multiplying the bottom row of the extended $p$-map for $\w^+$ by the inverse of the $p$-map matrix $(p_{ij})$ of $\w$. Here, as above, all equalities are understood in the locus $\{A_p^+\neq0\}\sse X(\beta_+)$.

%%%%%%%%%%%%%%%%%%%%%%%%%%%%%%%%%%%%%%%%%%%%%
%%%%%%%%%%%%%%%%%%%%%%%%%%%%%%%%%%%%%%%%%%%%%
\subsection{The cluster schemes of a Demazure weave}\label{ssec:A_and_X_schemes_weave} Building on \cite{FG2006,FG09}, the $\SA$-scheme and $\SX$-schemes are constructed from seed data in \cite[Section 2]{GHK15}. Given a Demazure weave $\w:\beta\to\delta(\beta)$, we denote by $\SA_\beta$ and $\SX_\beta$ the $\SA$-scheme and $\SX$-schemes associated to the seed data associated to $\w$. This seed data is typically encoded by the intersection quiver $Q_\w$ of Lusztig cycles. Intuitively, these schemes are obtained by gluing tori, see e.g.~\cite[Appendix A]{GHKK18}, indexed by the cluster $\SA$-seeds and the cluster $\SX$-seeds respectively:

$$\displaystyle\SA_\beta:=\bigcup_{{\bf s}\in S(\SA,\w)} T_{{\bf s}}^\SA,\qquad\qquad\displaystyle\SX_\beta:=\bigcup_{{\bf s}\in S(\SX,\w)} T^\SX_{{\bf s}}$$

\noindent where we have denoted by $S(\SA,\w)$ and $S(\SX,\w)$ the sets of all seeds obtained from the seed associated to $\w$ via a (possibly empty) sequence of cluster $\SA$-mutations and $\SX$-mutations, respectively. We refer to \cite[Prop.~2.4]{GHK15} for details on the gluing. As in \cite[Remark 2.7.$(iii)$]{GHK15}, there is a map $p:\SA_\beta\lr\SX_\beta$ which is compatible with the local maps between the respective tori. For the cluster structures associated to Demazure weaves, it follows from \cite[Section 8]{CGGLSS25} that there are scheme isomorphisms
\begin{equation}\label{eq:scheme_iso}
\Spec \Gamma(\SA_\beta,\mathcal{O}_{\SA_\beta}) \cong X(\beta),\qquad \Spec \Gamma(\SX_\beta,\mathcal{O}_{\SX_\beta}) \cong X(\beta).
\end{equation}
\noindent Thus, at the same time, the braid variety $X(\beta)$ can be considered as the affinization of the cluster $\SA$-scheme $\SA_\beta$, or as the affinization of the cluster $\SX$-scheme $\SX_\beta$. In the former viewpoint, the emphasis is on the fact that the coordinate ring $\kk[X(\beta)]$ is the cluster algebra $\SA(Q_\w)$, which in this case happens to be isomorphic to the upper cluster algebra $\Gamma(\SA_\beta,\mathcal{O}_{\SA_\beta})$. In the latter sense, the emphasis is on the fact that $X(\beta)$ admits a cluster Poisson structure, i.e.~a Poisson bracket which is given on each seed by the given skew-symmetric bilinear form. When we want to emphasize $X(\beta)$ as the affinization of the cluster $\SA$-scheme $\SA_\beta$, respectively the cluster $\SX$-scheme $\SX_\beta$, we will denote it by $\SA(\beta)$, respectively $\SX(\beta)$.\\

\begin{example}\label{ex:X_in_terms_of_A} For instance, given the map $p:\SA(\beta)\lr\SX(\beta)$, identity \eqref{eq:X_vars_via_A_vars} should be understood as expressing the pull-back of the $\SX$-variables under $p$ and expressing the resulting rational function in terms of the cluster $\SA$-variables:
$$p_\w^*(X^\w_{v_i}) := \prod_{v_j\in\w_3} (A^\w_{v_j}) ^{p_{ij}},$$
where $p_\w:T_\w^\SA\lr T_\w^\SX$ is the restriction of $p$ to the corresponding seed torus.\sbs
\end{example}

More generally, to ease notation, if we are discussing cluster $\SA$-variables, then $X(\beta)$ will implicitly mean $\SA(\beta)$, and if we discuss cluster $\SX$-variables, then $X(\beta)$ will refer to $\SX(\beta)$. If we discuss both, we will explicitly use $\SA(\beta)$ and $\SX(\beta)$. 

\begin{remark}\label{rmk:embedding_not_iso}
By \eqref{eq:scheme_iso}, there are natural embeddings
\begin{equation}\label{eq:embedding_not_iso}
\SA_\beta\hookrightarrow \SA(\beta),\qquad \SX_\beta\hookrightarrow \SX(\beta),
\end{equation}
but they are {\bf not} typically isomorphisms. Indeed, there are points in $X(\beta)$ that are not contained in any cluster $\SA$-seed or $\SX$-seed torus for many interesting positive braid words $\beta$. For instance, consider $\beta=\s_1^3$, then the point $(\borel,\borel_-,\borel,\borel_-)\in X(\beta)$ does not belong to any cluster seed, and so it is not in the image of either map in \eqref{eq:embedding_not_iso}. See \cite{CGSS_clusterdeeploci} for further discussions on this matter.\sbs
\end{remark}

In addition to the $\SA$ versus $\SX$ distinction, we recall that a Lie group $\G$ has been fixed in order to be able to consider the braid variety $X(\beta)$. Let us denote the braid variety $X(\beta)$ by $X_\G(\beta)$ when we want to emphasize the choice of $\G$. Following the convention above, $\SA_\G(\beta)$ denotes $X_\G(\beta)$ when seen as the affinization of the cluster $\SA$-scheme, and $\SX_\G(\beta)$ denotes $X_\G(\beta)$ when seen as the affinization of the cluster $\SX$-scheme. The notational emphasis on $\G$ is relevant in the context of cluster ensembles and Fock-Goncharov duality, cf.~\cite[Sections 1.2.9 \& 1.2.10]{FG09}, \cite[Definition A.4]{GHKK18} and \cite[Section 7]{GHK15}. Specifically, given $\G$, we denote by $\G^\vee$ the Langlands dual group, see e.g.~\cite[Appendix A]{GHKK18} or \cite[Section 1.17]{FG2006}. For instance if $\G=\SL_n$ then $\G^\vee=\operatorname{PSL}_n$. Then \cite[Theorem 8.8]{CGGLSS25} proves the following duality result on braid varieties:

\begin{theorem}[Fock-Goncharov duality for braid varieties]\label{thm:fock_goncharov_duality} Let $\beta$ be a positive braid word, $\G$ a simply-laced simple algebraic Lie group and $\G^\vee$ its Langlands dual. Then both
\begin{equation}\label{eq:cluster_ensembles}
(\SA_\G(\beta),\SX_{\G^\vee}(\beta))\qquad\qquad (\SA_{\G^\vee}(\beta),\SX_{\G}(\beta))
\end{equation}
are cluster ensembles, and Fock-Goncharov cluster duality holds for each of them.\sbs
\end{theorem}

\noindent For instance, \Cref{thm:fock_goncharov_duality} implies that the cluster algebra $\kk[\SA_\G(\beta)]$ admits a canonical basis of theta
functions parameterized by the integral tropical points of the dual braid variety $\SX_{\G^\vee}(\beta)$. The key aspect of the proof of \Cref{thm:fock_goncharov_duality} is showing that there exists a Donaldson-Thomas transformation\footnote{Combinatorially, that there exists a reddening sequence for the quiver, cf.~\cite{Keller2017_DT,Keller12_DT2}.} and that the exchange matrices from Demazure weaves are full rank (for the extended exchange matrix with mutable columns). In fact, there is a rather explicit map $D_\beta:X(\beta)\lr X(\beta)$, generalizing the usual twist automorphisms, that realizes the Donaldson-Thomas transformation, cf.~\cite[Theorem 8.7]{CGGLSS25}.

%%%%%%%%%%%%%%%%%%%%%%%%%%%%%%%%%%%%%%%%%%%%%
%%%%%%%%%%%%%%%%%%%%%%%%%%%%%%%%%%%%%%%%%%%%%
%%%%%%%%%%%%%%%%%%%%%%%%%%%%%%%%%%%%%%%%%%%%%
%%%%%%%%%%%%%%%%%%%%%%%%%%%%%%%%%%%%%%%%%%%%%
\subsection{Brick compactifications of a braid variety}\label{ssec:brick_compactifications} Let $\beta=\s_{i_1}\cdots\s_{i_\ell}$ be a positive braid word and $X(\beta)$ its braid variety. Following \cite[Section 4]{CGGS2}, we now construct a smooth projective compactification $Y(\beta)$ of $X(\beta)$. As we will explain, such compactification $Y(\beta)$ has a number of desired properties both from the combinatorial viewpoint and as it relates to aspects of mirror symmetry, cf.~\Cref{thm:superpotential_Reeb_coordinates}. Importantly, $Y(\beta)$ depends on the choice of positive braid word $\beta$ and not just on its underlying braid $[\beta]$. Therefore, since the isomorphism type of $X(\beta)$ only depends on the underlying braid $[\beta]$, we can construct several well-behaved compactifications of $X(\beta)$ by using different positive braid words for $[\beta]$.\\

\noindent Given a positive braid word $\beta$, the intuition for compactifying $X(\beta)$ is simple: we allow pairs of consecutive flags in \Cref{def:braid_variety} to be equal, not just in a certain relative position. This intuition is formalized as follows.

%%%%%%%%%%%%%%%%%%%%%%%%%%%%%%%%%%%%%%%%%%%%%
%%%%%%%%%%%%%%%%%%%%%%%%%%%%%%%%%%%%%%%%%%%%%
\subsubsection{The brick compactification $Y(\beta)$}\label{sssec:brick_compactification} Let \(P_i\) denote the minimal parabolic subgroup associated with
\(s_i\) which contains the fixed Borel $\borel\sse\G$. For instance, if $\G=\SL_n$ and $\borel$ is the subgroup of invertible upper-triangular matrices, then
$$P_i\cong\{(M_{jk})\in\SL_n: M_{jk}=0\mbox{ if }k<j\mbox{ and }(j,k)\neq(i+1,i)\}.$$
That is, the subgroup of invertible upper-triangular matrices where, in addition, the below-diagonal entry $(i+1,i)$ is allowed to be non-zero.\\

By definition, the closed Bott--Samelson variety associated to the braid word $\beta=\s_{i_1}\ldots\s_{i_{\ell}}$ is
\begin{equation}\label{eq:BSquotient}
 \BS(\beta)=(P_{i_1}\times_\borel P_{i_2}\times_\borel\cdots\times_\borel P_{i_{\ell}})/\borel^{\ell},
\end{equation}
where $\borel^{\ell+1}$ acts on the right by
$$(p_1,\ldots,p_{\ell})\cdot (b_1,\ldots,b_{\ell}):=(p_1b_1,b_1^{-1}p_2b_2,\ldots,b_{\ell-1}^{-1}p_{\ell} b_{\ell}).$$
That is, $\BS(\beta)$ parametrizes the following tuples of flags:
\begin{equation}\label{eq:BSgallery}
 \BS(\beta)=\Bigl\{(\borel_1,\ldots,\borel_{\ell+1}):
 \borel_1=\borel,\quad \borel_{k} \buildrel s_{i_{k}}\mbox{\tiny or }e \over \longrightarrow \borel_{k}\}.
\end{equation}
For any braid word $\beta$, $\BS(\beta)$ is a smooth projective variety\footnote{In fact, it is an iterated $\PP^1$-bundle, and thus smooth and projective.} of dimension $\dim_\kk\BS(\beta)=\ell$. Consider the endpoint map
\begin{equation}\label{eq:endpoint}
  m_\beta:\BS(\beta)\longrightarrow \G/\borel,
  \qquad (\borel_1,\ldots,\borel_{\ell+1})\longmapsto \borel_{\ell+1}.
\end{equation}
For instance, the image of $m_\beta$ is the classical Schubert variety $\overline{B\dot uB/B}$ with $u=\delta(\beta)$.

\begin{definition}[The compactification $Y(\beta)$]\label{def:brick}
Let $\beta=\s_{i_1}\cdots\s_{i_\ell}$ be a positive braid word. By definition, the brick compactification $Y(\beta)$ associated to $\beta$ is
\begin{equation}\label{eq:brickdef}
  Y(\beta):=m_\beta^{-1}(\delta(\beta)\borel).
\end{equation}
where $\delta(\beta)\in W(\G)$ has been lifted to $N_\G(\T)$. Equivalently,
\begin{equation}\label{eq:brickgallery}
Y(\beta):=\Bigl\{(\borel_1,\ldots,\borel_{\ell+1}):
 \borel_1=\borel, \borel_{k} \buildrel s_{i_{k}}\mbox{\tiny or }e \over \longrightarrow \borel_{k+1},\borel_{\ell+1}=\delta(\beta)\borel\}.
\end{equation}
\sbs
\end{definition}

These varieties were introduced and studied in \cite{Escobar}. It is proven in \cite[Theorem 20]{Escobar} that $Y(\beta)$ is a smooth irreducible projective variety of dimension $\dim_\kk Y(\beta)=\ell(\beta)-\ell(\delta(\beta))$, which coincides with the dimension of $X(\beta)$. By comparing \Cref{def:brick}, especially \eqref{eq:brickgallery}, with \Cref{def:braid_variety} it thus follows that $X(\beta)\sse Y(\beta)$ is an open (and dense) subvariety.

%%%%%%%%%%%%%%%%%%%%%%%%%%%%%%%%%%%%%%%%%%%%%
%%%%%%%%%%%%%%%%%%%%%%%%%%%%%%%%%%%%%%%%%%%%%
\subsubsection{The boundary divisors}\label{sssec:dual_boundary_complex}
The complement $Y(\beta)\setminus X(\beta)$ is a union of irreducible boundary divisors $\{D_i\}$. Specifically, we can consider
\begin{equation}\label{eq:Djdef}
  D_j:=\{(\borel_1,\ldots,\borel_{\ell+1})\in Y(\beta):\borel_{j}=\borel_{j+1}\},\quad j\in[1,\ell].
\end{equation}
Let us denote by $\beta\setminus \s_{i_j}$ the positive braid word resulting from deleting the $j$th crossing $\s_{i_j}$ from $\beta$. More generally, if $S$ is a subset of crossings of $\beta$, we denote by $\beta\setminus S$ the positive braid word resulting from deleting the crossings in $S$ from $\beta$. The following holds:

\begin{lemma}[Boundary divisors]\label{lem:boundary_div} Let $\beta$ be a positive braid word, $X(\beta)$ its braid variety and $Y(\beta)$ its brick compactification. Then:
\begin{enumerate}
    \item $D_j\ne\emptyset
  \quad\Longleftrightarrow\quad
  \delta(\beta\setminus \s_{i_j})=\delta(\beta)$\\

  \item If $\delta(\beta\setminus \s_{i_j})=\delta(\beta)$, then $D_j\sse Y(\beta)$ is a smooth irreducible divisor and $D_j\cong Y(\beta\setminus \s_{i_j})$.\\

  \item More generally, $\displaystyle\bigcap_{\s_{i_j}\in S}D_j\ne\emptyset\iff \delta(\beta\setminus S)=\delta(\beta)$ and, in such a case,
  \begin{equation}\label{eq:multiisomorphism}
 \bigcap_{\s_{i_j}\in S}D_j\cong Y(\beta\setminus S).
\end{equation}
In addition, such intersection in \eqref{eq:multiisomorphism} has codimension $|S|$ in $Y(\beta)$.\\
 \item The union of divisors
 $$\bigcup_{j}D_j\sse Y(\beta),$$
 where the union is over indices $j\in[1,\ell]$ satisfying $\delta(\beta\setminus \s_{i_j})=\delta(\beta)$, is a simple normal crossings divisor in $Y(\beta)$.\\

  \item The dual boundary complex of $\{D_j\}$ is isomorphic to the subword complex of $(\beta,\delta(\beta))$. In addition, it is homotopy equivalent to a sphere.
  \sbs
\end{enumerate}
\end{lemma}
Parts \Cref{lem:boundary_div}.(1) and (2) follow from direct computation, \Cref{lem:boundary_div}.(3) follows from iterating \Cref{lem:boundary_div}.(2). \Cref{lem:boundary_div}.(4) follows from \Cref{lem:boundary_div}.(3), as the divisors intersect transversely, since all the intersections are smooth, cf.~\cite[Theorem 4.11]{CGGS2} for more details. The subword complexes mentioned in \Cref{lem:boundary_div}.(5) were studied in \cite{KnutsonMiller_04} and sphericity follows from \cite[Corollary 3.8]{KnutsonMiller_04}, see also \cite[Section 4]{CGGS2}.\footnote{The assertion that the dual boundary complex is homotopy equivalent to a sphere is an instance of the geometric $P=W$ conjecture.}\\

\noindent Let $D_\beta:=\sum D_j$ denote the linear sum of the boundary divisors of the brick compactification $Y(\beta)$ of $X(\beta)$. Beyond \Cref{lem:boundary_div}, the following property of $D_\beta$ is significant from the viewpoint of mirror symmetry and the study of cluster algebras via algebraic geometry:

\begin{proposition}\label{prop:logCY}
Let $\beta$ be a positive braid word, $X(\beta)$ its braid variety, $Y(\beta)$ its brick compactification and $D_\beta$ the linear sum of its compactifying divisors. Then:
\begin{enumerate}
    \item $(Y(\beta),D_\beta)$ is a log Calabi-Yau pair.
    \item $(Y(\beta),D_\beta)$ has maximal boundary and $D_\beta$ is the support of an ample divisor.\sbs
\end{enumerate}

\end{proposition}

\noindent Equivalently, after base change\footnote{If we are to remain over $\Z$, then the resulting object is a relative regular form.} to $\C$, \Cref{prop:logCY}.(1) states that there exists a nowhere vanishing holomorphic top form on $X(\beta)$ which acquires only simple poles along $D_\beta$ when it compactifies to $Y(\beta)$, see \cite{HackingKeel_ICM18} for details on log Calabi-Yau pairs as they relate to cluster structures. \Cref{prop:logCY} follows from studying the divisors and canonical bundle of $\BS(\beta)$, see e.g.~\cite{Anderson,LauTho} or \cite[Lemma 4.9]{CasalsWeng22}. For braid varieties $X(\beta)$ such that the closed 2-form $\omega_\beta$ in \Cref{sssec:symplectic_form_coordinates} is non-degenerate, we can directly take the top power of $\omega_\beta$.\\

\noindent The maximal-boundary statement in \Cref{prop:logCY}.(2), i.e.~that there is a 0-dimensional stratum in the intersection complex of $D_\beta$, follows from \Cref{lem:boundary_div}. Note that \Cref{prop:logCY}.(2) only states that $D_\beta$ is the support of an ample divisor, but $D_\beta$ itself might not be ample. Indeed, by \Cref{prop:logCY}.(1) $D_\beta$ is itself ample iff $Y(\beta)$ is a Fano variety, which is not always the case.%, e.g.~for $\beta=c\cdot S_c(w_0)$ where $c$ represents the Coxeter element and $S_c(w_0)$ is the $c$-sorting reduced word for $w_0$, $Y(\beta)$ is a toric variety obtained from the associahedron, as in \cite[Corollary 2.10]{Escobar}, and it is not Fano for $\G=\SL_n$ if $n\geq3$.

%%%%%%%%%%%%%%%%%%%%%%%%%%%%%%%%%%%%%%%%%%%%%
%%%%%%%%%%%%%%%%%%%%%%%%%%%%%%%%%%%%%%%%%%%%%
\subsubsection{Two applications of brick compactifications}\label{sssec:applications_compactifications}

First, the cohomology of the variety $X(\beta)$ can be computed by using the fact that the cohomology of $Y(\beta)$ is rather well-understood and applying the inclusion-exclusion principle. Specifically, consider the $\T$-action on $X(\beta)$ introduced in \cref{sssec:braid_variety_further_symmetries}, cf.~\eqref{eq:torus_action}, and let $H_{c,T}^*(X(\beta),\Z)$ denote the compactly supported $\T$-equivariant cohomology of $X(\beta)$. Since $X(\beta)$ is a smooth affine variety, $H_{c,T}^*(X(\beta),\Z)$ is equipped with a weight filtration (after base change from $\Z$ to $\C$) and we denote by $\operatorname{gr}^W H_{c,T}^*(X(\beta),\Z)$ the associated graded quotients of this weight filtration.\\

A neat result about braid varieties is that such graded pieces of the equivariant cohomology of $X(\beta)$ recover a link invariant if $\G=\SL_n$. Indeed, let $\operatorname{HHH}$ be the triply-graded Khovanov-Rozansky link homology, as in \cite{KhS,KR}, and we denote by $\operatorname{HHH}(\beta w_0)$ the triply-graded Khovanov-Rozansky link homology of the smooth link obtained by the closure of $\beta w_0$, where $w_0$ denotes a braid lift of a reduced expression for $w_0$. Recall that one of the three degrees is the $a$-degree, which ranges from $0$ to the number of strands of $\beta$. Let us denote by $\operatorname{HHH}^{a=n}$ the highest $a$-degree piece and we use $\ateb$ to denote $\beta$ read right-to-left. The following result is then proven in \cite[Theorem 4.22]{CGGS2}:

\begin{theorem}\label{thm:cohomology_and_HHH}
Let $\beta$ be a positive $n$-stranded braid word with $\delta(\beta)=w_0$. Then
\begin{equation}\label{eq:cohomology_and_HHH}
\boxed{\operatorname{gr}^W H_{c,T}^*(X(\beta),\Z)\cong \operatorname{HHH}^{a=n}(\ateb w_0)}
\end{equation}
\sbs
\end{theorem}

\noindent \Cref{thm:cohomology_and_HHH} is a cohomological enhancement of the fact that the ruling polynomial of a Legendrian link appears as the appropriate $a$-coefficient  of the HOMFLY polynomial of the underlying smooth link, see e.g.~\cite{Rutherford06_HOMFLY,Rutherford11_HOMFLY}. It can also be understood as a cohomological enhancement in characteristic zero of certain identities between finite points counts, cf.~\cite{trinh2021heckecategoryunipotentlocus,trinh2022unipotentelementstwistinglink}.\\

Second, it is possible to use $Y(\beta)$ as a compactification of $X(\beta)$ in order to establish a form of mirror symmetry for the appropriate cluster $\SA$ and $\SX$-schemes associated to $\beta$. Specifically, in the context of Fock-Goncharov duality, as in \Cref{thm:fock_goncharov_duality}, the cluster ensemble
$$(\SA_\G(\beta),\SX_{\G^\vee}(\beta))$$
can be refined by compactifying $\SX_{\G^\vee}(\beta)$ and, dually, adding a potential $W_\beta:\SA_\G(\beta)\lr\mathbb{A}^1_\kk$ to $\SA_\G(\beta)$. This turns $(\SA_\G(\beta),W_\beta)$ into a Landau-Ginzburg model from which it is possible to compute certain invariants. These potentials $W_\beta$ have clean expressions in terms of Reeb chord asymptotics, as we explain in \Cref{thm:superpotential_Reeb_coordinates} below.

\begin{remark} We use the notation $W_\beta$ because the potential depends on the positive braid word $\beta$, as does the compactification $Y(\beta)$ of $\SX_{\G^\vee}(\beta)$. Thus, as above, one can change the positive braid word $\beta$ by braid moves and obtain different Landau-Ginzburg models for the same variety $\SA_\G(\beta)$.\sbs
\end{remark}